\documentclass[conference]{IEEEtran}
\IEEEoverridecommandlockouts
% The preceding line is only needed to identify funding in the first footnote. If that is unneeded, please comment it out.
\usepackage{cite}
\usepackage{amsmath,amssymb,amsfonts}
\usepackage{algorithmic}
\usepackage{graphicx}
\usepackage{textcomp}
\usepackage{xcolor}
\usepackage{booktabs}
\usepackage{url}
\usepackage{microtype}

\def\BibTeX{{\rm B\kern-.05em{\sc i\kern-.025em b}\kern-.08em
    T\kern-.1667em\lower.7ex\hbox{E}\kern-.125emX}}

\begin{document}

\title{EvidenceGain-X: Counterfactual, Self-Verifying Next-Best Evidence Acquisition for Crop Disease Diagnosis}

\author{
    \IEEEauthorblockN{Rishi Choudhary\IEEEauthorrefmark{1}, Kumar Sanib\IEEEauthorrefmark{1}, Amrit Raj\IEEEauthorrefmark{1}, Nikhil Yadav\IEEEauthorrefmark{1}, Amit Venkat\IEEEauthorrefmark{1}, Nihal Kakollu\IEEEauthorrefmark{1}, and Devulapalli Sudheer\IEEEauthorrefmark{2}}
    \IEEEauthorblockA{Department of Computer Science and Engineering, School of Computer Science and Engineering (SCOPE)\\
    VIT-AP University, Amaravati, Andhra Pradesh 522237, India\\
    Emails: \IEEEauthorrefmark{1}\{rishi.24bce7490, kumar.24bca7144, amrit.24bce7856, nikhil.24bca7057, amit.24bca7174, nihal.24bca7283\}@vitapstudent.ac.in\\
    \IEEEauthorrefmark{2}sudheer.d@vitap.ac.in}
}

\maketitle

\begin{abstract}
Deployed computer vision systems for agricultural plant pathology overwhelmingly operate under a passive, single-shot classification regime: an end-user submits a single foliar photograph, and a convolutional neural network outputs an uncalibrated point prediction. While this paradigm yields benchmark accuracies exceeding 95\% on curated datasets, it fails in practical field environments due to four persistent vulnerabilities: (i) acute morphological look-alike ambiguity along early vegetative decision boundaries (\textit{Alternaria solani} vs. \textit{Septoria lycopersici}), (ii) physical occlusion of pathognomonic discriminators (such as abaxial downy sporulation in \textit{Phytophthora infestans}), (iii) unconstrained Softmax confidence hallucination on optically degraded specimens, and (iv) Bayesian posterior corruption caused by inadvertent user uploads of healthy or foreign vegetation.

To resolve these failure modes, this paper presents EvidenceGain-X, an autonomous, closed-loop diagnostic framework comprising nine functionally distinct modules that reformulates plant pathology from passive inference into an active, self-verifying clinical dialogue. EvidenceGain-X quantifies its predictive uncertainty through temperature-scaled Shannon entropy ($T = 1.8928$, optimized via L-BFGS to reduce Expected Calibration Error by 78.2\%), simulates counterfactual posterior belief trajectories under competing diagnostic hypotheses using an empirical discriminative affinity matrix, and determines the optimal next-best evidence query ($e^*$) by maximizing a cost-aware Net Policy Utility function across five multimodal modalities. Crucially, the architecture incorporates two cyber-physical safety shields: Module 7 Optical Quality Gating (discrete Laplacian variance $\sigma^2 \ge 80.0$ and luminance bounds $[30, 230]$\,cd/m$^2$) enforcing deterministic blur/glare interception on both initial and follow-up inputs, and Module 8 Contradiction \& Specimen Mismatch Guard (directional cosine divergence $\ge 0.40$ and out-of-distribution prior-evidence bounds), which quarantines incongruent foliage.

Evaluated on 24 multi-view diagnostic episodes from the PlantVillage tomato benchmark under leakage-free plant-grouped partitioning (70/15/15), EvidenceGain-X elevates diagnostic accuracy from 87.50\% (passive single-shot) to 91.67\% and attains a Macro F1 score of 0.9188 in an average of only 0.88 sequential observation steps, while achieving the highest Shannon entropy reduction (0.359 bits) among all evaluated policies. Under adversarial and noisy conditions, EvidenceGain-X achieves a 100\% rejection rate against motion-blurred, underexposed, and cross-specimen contaminant inputs. The entire pipeline operates with an edge CPU latency under 30\,ms and a memory footprint of 44.5\,MB, establishing an operational standard for trustworthy edge AI in digital agriculture.
\end{abstract}

\begin{IEEEkeywords}
Active sequential vision, Bayesian belief updating, counterfactual planning, crop disease diagnosis, foliar pathology, optical quality gating, precision agriculture, self-verifying AI, temperature scaling.
\end{IEEEkeywords}

\section{Introduction}

\subsection{Agronomic Context and the Crisis of Passive Vision AI}
Foliar crop pathologies account for an estimated 20\% to 40\% reduction in global crop yields annually, imperiling food security and economic stability for over 500 million smallholder farming households across South Asia and sub-Saharan Africa \cite{hughes2015, mohanty2016}. Among staple horticultural commodities, tomato (\textit{Solanum lycopersicum}) production exceeds 180 million metric tonnes annually; however, the crop is exceptionally susceptible to aggressive fungal and oomycete pathogens, primarily Early Blight (\textit{Alternaria solani}), Late Blight (\textit{Phytophthora infestans}), and Septoria Leaf Spot (\textit{Septoria lycopersici}) \cite{agrios2005}. In particular, \textit{Phytophthora infestans} exhibits explosive infection dynamics, capable of transforming a healthy canopy into completely necrotic tissue within 7 to 10 days under humid, temperate microclimates. The asymmetric economic hazard of diagnostic error---where a false negative for Late Blight results in catastrophic harvest liquidation---demands diagnostic systems that exhibit uncompromising clinical reliability \cite{ferentinos2018}.

Over the past decade, deep learning architectures---most notably deep Convolutional Neural Networks (CNNs) and Vision Transformers---have achieved reported classification accuracies exceeding 98\% on curated agricultural benchmark repositories like PlantVillage \cite{hughes2015, mohanty2016}. Despite these high benchmark scores, practical field deployment of mobile agricultural diagnostic tools remains plagued by severe misdiagnoses. Our empirical investigation reveals four structural failure modes intrinsic to the conventional \textbf{passive single-shot paradigm}:

\begin{enumerate}
    \item \textbf{Morphological Look-Alike Ambiguity:} In early infection stages, distinct fungal taxa converge morphologically. \textit{Alternaria solani} and \textit{Septoria lycopersici} both manifest as small, circular, brown necrotic flecks encircled by chlorotic yellow halos. In top-down macroscopic photography, spatial sensor resolution is fundamentally insufficient to resolve microscopic pycnidia (black fungal fruiting bodies diagnostic of Septoria) from immature concentric target rings. Standard classifiers force an uncalibrated, high-entropy guess between bimodal posteriors ($P \approx 50\%$ vs. $50\%$).
    \item \textbf{Observational Occlusion of Diagnostic Discriminators:} Plant pathologies are intrinsically three-dimensional, anatomical phenomena. Pathognomonic symptoms frequently develop on leaf structures concealed from conventional overhead viewpoints. Specifically, the defining morphological proof of Late Blight is delicate, white downy sporulation located exclusively on the \textbf{abaxial (underside) leaf lamina}, or dark petiole lesions along the stem. A single macroscopic photograph cannot inspect what is physically occluded.
    \item \textbf{Softmax Hallucination on Degraded Imagery:} Standard neural network heads apply the Softmax activation function over final classification logits $\mathbf{z}$:
    \begin{equation}
        \sigma(\mathbf{z})_i = \frac{e^{z_i}}{\sum_{j=1}^K e^{z_j}}
    \end{equation}
    Because Softmax mathematically normalizes logits to sum to unity ($\sum \sigma(\mathbf{z})_i = 1$), it lacks an explicit ``uncertain'' or ``corrupted input'' state. When a farmer captures an out-of-focus, motion-blurred, or severely underexposed photograph, the convolutional feature extractor maps low-frequency green blur to arbitrary high-confidence predictions (e.g., assigning 89.5\% probability of Late Blight to empty blur), creating severe safety risks.
    \item \textbf{Cross-Specimen Contamination \& Contradiction:} In sequential multimodal inquiry, non-expert users frequently photograph an adjacent asymptomatic leaf, capture a weed, or ingest an unrelated specimen. Conventional feature-averaging or Bayesian fusion architectures blindly absorb contradictory inputs, corrupting prior diagnostic beliefs and destabilizing the clinical decision path.
\end{enumerate}

\subsection{The EvidenceGain-X Paradigm: Active Clinical Vision}
To overcome these limitations, this paper introduces \textbf{EvidenceGain-X}, an autonomous diagnostic system modeled upon human clinical reasoning. When an expert plant pathologist encounters an ambiguous foliar presentation, they do not commit to an immediate guess. Instead, they quantify diagnostic uncertainty, formulate competing differential hypotheses ($h_1$ vs. $h_2$), simulate what information prospective physical examinations will yield, and instruct the farmer to perform a specific targeted maneuver (e.g., \textit{``turn the leaf over to inspect the underside''}). Furthermore, incoming evidence is rigorously vetted for photographic clarity and specimen congruence before clinical belief is updated.

EvidenceGain-X formalizes this clinical methodology into a nine-module closed-loop pipeline. By integrating temperature-scaled uncertainty gating, counterfactual trajectory simulation, multi-objective utility optimization, optical blur gating, and directional contradiction verification, EvidenceGain-X transforms crop pathology from passive single-shot inference into an active, self-correcting clinical dialogue.

\subsection{Key Contributions}
The primary contributions of this work are summarized as follows:
\begin{itemize}
    \item \textbf{Counterfactual Evidence Planning (Module 5):} A goal-directed planning engine that models hypothetical posterior trajectories under competing disease hypotheses, employing a domain-specific discriminative affinity matrix ($\mathbf{A} \in \mathbb{R}^{M \times K \times K}$) calibrated from agronomic literature to compute expected information gain without requiring intractable Monte Carlo approximations.
    \item \textbf{Cost-Aware Multi-Factor Utility Function (Module 6):} A multi-objective query selection policy combining five Pareto hyperparameters ($\alpha_{\text{sep}} = 1.0$, $\beta_{\text{unc}} = 0.8$, $\gamma_{\text{diag}} = 1.2$, $\lambda_{\text{cost}} = 0.3$, $\mu_{\text{risk}} = 0.5$) that balances expected entropy reduction against physical farmer acquisition effort and operational failure risk.
    \item \textbf{Dual Cyber-Physical Safety Shields (Modules 7 and 8):}
    \begin{itemize}
        \item \textit{Module 7 (Optical Quality Gating):} Deterministic pre-screening enforcing Laplacian edge variance ($\sigma^2 \ge 80.0$) and luminance bounds ($[30, 230]$\,cd/m$^2$) on \textbf{both initial specimen ingestion ($X_0$) and sequential evidence}, eliminating Softmax hallucination on degraded imagery.
        \item \textit{Module 8 (Contradiction \& Specimen Mismatch Guard):} Dual-criterion verification that detects directional cosine divergence ($S_{\text{contra}} \ge 0.40$) and foreign pathology intrusion ($\theta_{\text{ev}} \notin \{h_1, h_2\}$ with prior $< 0.05$ and evidence $\ge 0.35$), quarantining corrupted uploads and preserving prior diagnostic states.
    \end{itemize}
    \item \textbf{Leakage-Safe Validation \& Open-Source Implementation:} Complete empirical validation across 24 PlantVillage episodes partitioned under physical plant-level isolation (zero inter-partition leaf leakage), accompanied by an open-source, edge-deployable implementation with a 44.5\,MB footprint and sub-30\,ms CPU inference latency.
\end{itemize}

\section{Related Work}

\subsection{Deep Learning for Plant Disease Classification}
The PlantVillage dataset, introduced by Hughes and Salath\'e \cite{hughes2015} and expanded by Mohanty \textit{et al.} \cite{mohanty2016}, established the first large-scale benchmark for automated foliar pathology identification, comprising over 54,000 images across 26 diseases. A succession of architectures---from AlexNet and GoogLeNet through ResNet \cite{he2016}, DenseNet, and Vision Transformers---have achieved reported accuracies exceeding 99\% on random hold-out splits. However, Ferentinos \cite{ferentinos2018} demonstrated that these figures are substantially inflated by data leakage: photographs of the same physical leaf appear in both training and test partitions under random splitting. Plant-grouped partitioning, as employed in this work, eliminates this leakage and typically reduces reported accuracies by 5 to 15 percentage points.

Beyond the leakage issue, existing approaches treat plant disease diagnosis as a passive, single-shot classification problem, assuming one image contains sufficient information for diagnosis. Ramcharan \textit{et al.} \cite{ramcharan2017} demonstrated that field-collected cassava images produced significantly lower accuracies than laboratory datasets, underscoring the gap between benchmark performance and real-world deployment. EvidenceGain-X addresses this fundamental limitation by transforming the single-shot paradigm into a sequential evidence-acquisition framework where the system actively requests the most diagnostically informative observations.

\subsection{Uncertainty Quantification and Calibration}
Guo \textit{et al.} \cite{guo2017} demonstrated that modern deep networks trained with cross-entropy loss produce systematically overconfident softmax distributions---a phenomenon termed miscalibration. Temperature scaling, a post-hoc method in which logits are divided by a learned scalar $T$ before softmax, substantially reduces Expected Calibration Error (ECE) without altering discriminative accuracy. In EvidenceGain-X, calibrated uncertainty is not merely desirable but prerequisite: the stopping criterion, query selection, and information gain computation all depend on Shannon entropy being a faithful proxy for genuine diagnostic ambiguity.

Alternative uncertainty methods include Monte Carlo Dropout \cite{gal2016}, which requires multiple stochastic forward passes, and deep ensembles \cite{lakshminarayanan2017}, which multiply model size linearly. Temperature scaling adds zero inference-time computation, making it the only calibration method strictly compatible with low-power agricultural edge-deployment constraints.

\subsection{Active Learning and Sequential Experimental Design}
Pool-based active learning \cite{settles2009} selects unlabelled instances for annotation by a human oracle to maximize model improvement per annotation cost. The setting in EvidenceGain-X differs in a crucial respect: we acquire complementary physical observations of a single diagnostic case to refine a running posterior distribution, rather than labels to retrain a model. This distinction aligns the framework with Bayesian optimal experimental design \cite{chaloner1995, lindley1956}, which formalizes sequential evidence acquisition as maximizing expected information gain---the expected reduction in Shannon entropy of the posterior.

For deep learning models, computing exact expected information gain is generally intractable. EvidenceGain-X decomposes the computation into two tractable components: a domain-specific discriminative affinity matrix that encodes prior agronomic knowledge about which evidence modalities separate which disease pairs, and a lightweight neural utility regressor (the UtilityScoringNetwork) trained on empirical gain targets from historical episodes. This decomposition enables real-time query selection on resource-constrained edge devices.

\section{The Proposed EvidenceGain-X Framework}

The complete EvidenceGain-X pipeline comprises nine functionally distinct modules arranged in a closed-loop architecture. An initial macroscopic leaf photograph enters the pipeline and is pre-screened for optical clarity (Module 7), after which visual features are extracted (Module 1) and a calibrated provisional diagnosis (Module 2) is generated. An ambiguity decision gate evaluates whether the diagnosis is sufficiently confident to terminate. If not, the system initiates a sequential inquiry loop involving disease-confusion analysis (Module 3), candidate evidence generation (Module 4), counterfactual evidence planning (Module 5), multi-factor utility scoring (Module 6), dual safety verification (Modules 7 and 8), and Bayesian belief updating (Module 9). The loop iterates until the stopping criterion is satisfied or the maximum evidence budget is exhausted.

\subsection{Module 1: Vision Backbone Encoder}
The initial leaf photograph $X_0 \in \mathbb{R}^{H \times W \times 3}$ is resized bicubically to $224 \times 224$ pixels and normalized using ImageNet channel statistics. Feature extraction uses an ImageNet-pretrained ResNet-18 architecture \cite{he2016}:
\begin{equation}
    \mathbf{h}_0 = f_{\phi}(X_0) \in \mathbb{R}^{512}
\end{equation}
A task-specific linear classification head $\mathbf{W}_{\text{cls}} \in \mathbb{R}^{4 \times 512}$ with bias $\mathbf{b}_{\text{cls}} \in \mathbb{R}^4$ maps the latent feature vector $\mathbf{h}_0$ to unnormalized diagnostic logits:
\begin{equation}
    \mathbf{z} = \mathbf{W}_{\text{cls}} \mathbf{h}_0 + \mathbf{b}_{\text{cls}}
\end{equation}
ResNet-18 was selected over heavier backbones (ResNet-50, ViT) based on empirical edge execution constraints: at 44.0\,MB checkpoint size and 12\,ms CPU forward-pass latency, it delivers real-time inference on edge processors while preventing overfitting on specialized agricultural subsets.

\subsection{Module 2: Calibrated Provisional Diagnosis}
To correct deep network overconfidence, raw classification logits are calibrated via post-hoc Platt Temperature Scaling \cite{guo2017}:
\begin{equation}
    P(\theta_c \mid X_0) = \frac{\exp(z_c / T)}{\sum_{j=1}^K \exp(z_j / T)}, \quad T = 1.8928
\end{equation}
where $K = 4$ denotes the target pathology taxonomy: $\Theta = \{\text{Early Blight}, \text{Late Blight}, \text{Septoria Leaf Spot}, \text{Healthy}\}$.

The temperature parameter $T^* = 1.8928$ is learned by minimizing negative log-likelihood over the uncorrupted validation partition using L-BFGS, keeping underlying backbone and classification head weights frozen. From the calibrated distribution, three diagnostic scalars are computed:
\begin{align}
    H(P) &= -\sum_{c=1}^K P(\theta_c) \log_2 P(\theta_c) \\
    h_1 &= \operatorname*{arg\,max}_{\theta_c \in \Theta} P(\theta_c), \quad h_2 = \operatorname*{arg\,max}_{\theta_c \in \Theta \setminus \{h_1\}} P(\theta_c) \\
    \Delta p &= P(h_1) - P(h_2)
\end{align}

\subsection{Ambiguity Decision Gate and Stopping Rules}
The decision gate evaluates two conditions jointly:
\begin{equation}
    \text{TERMINATE} \iff \left[ H(P) \le \tau_H \right] \;\wedge\; \left[ \Delta p \ge \tau_m \right]
\end{equation}
where $\tau_H = 0.35$\,bits and $\tau_m = 0.70$. If both conditions are satisfied simultaneously, the system commits to diagnosis $h_1$ at zero additional observation cost. If either condition is violated, the diagnosis is declared ambiguous, and active inquiry is initiated.

\subsection{Module 3: Disease-Confusion Analysis}
When the decision gate triggers active inquiry, Module 3 inspects probability mass concentrations across the rival pair $(h_1, h_2)$ to classify the specific clinical confusion typology:
\begin{itemize}
    \item \textit{Systemic Blight Confusion:} $P(\text{Early Blight}) + P(\text{Late Blight}) \ge 0.75$, where macroscopic foliar necrosis requires pathogen-specific etiology determination.
    \item \textit{Foliar Spot Confusion:} $P(\text{Early Blight}) + P(\text{Septoria}) \ge 0.75$, representing near-identical circular lesions requiring microscopic structural discrimination (pycnidia vs. concentric rings).
\end{itemize}

\subsection{Module 4: Candidate Evidence Pool and Cost Modeling}
EvidenceGain-X maintains an operational catalog of five candidate multimodal evidence modalities $\mathcal{E}$, parameterized by acquisition cost $C(e)$ and physical capture failure risk $R(e)$:
\begin{enumerate}
    \item \textit{Lesion Close-Up ($e_1$):} Macro zoom ($C = 1.0, R = 0.20$).
    \item \textit{Abaxial Leaf Underside ($e_2$):} Inverted leaf view revealing sporulation ($C = 1.2, R = 0.30$).
    \item \textit{Stem Petiole View ($e_3$):} Stem photograph inspecting cankers ($C = 1.5, R = 0.35$).
    \item \textit{Field Symptom Description ($e_4$):} Textual observation encoded via TextEvidenceEncoder ($C = 0.4, R = 0.10$).
    \item \textit{Microclimate Context ($e_5$):} Temperature, humidity, and rainfall context ($C = 0.4, R = 0.15$).
\end{enumerate}
Acquired modalities are removed dynamically from $\mathcal{E}$ to prevent redundant querying. A maximum horizon budget $B = 3$ limits total farmer burden.

\subsection{Module 5: Counterfactual Evidence Planning}
For each available clue $e \in \mathcal{E}$, the planner simulates two hypothetical posterior distributions: $P(h \mid e, h_1)$ (if evidence supports $h_1$) and $P(h \mid e, h_2)$ (if evidence supports $h_2$). Simulations are governed by an empirical Discriminative Affinity Matrix $\mathbf{A} \in \mathbb{R}^{|\mathcal{E}| \times K \times K}$, where entry $A[e, i, j] \in [0, 1]$ represents the domain-specific power of modality $e$ to separate pathology $i$ from pathology $j$. Under the counterfactual assumption that $e$ supports $h_1$, probability mass is transferred:
\begin{equation}
    \Delta P = s \cdot \kappa \cdot P(h_2), \quad s = A[e, h_1, h_2], \; \kappa = 0.7
\end{equation}
The divergence between competing simulated futures is measured via the symmetric Jensen-Shannon Divergence (JSD):
\begin{equation}
    D_{\text{JS}}(P_1 \parallel P_2) = \frac{1}{2} D_{\text{KL}}(P_1 \parallel M) + \frac{1}{2} D_{\text{KL}}(P_2 \parallel M)
\end{equation}
where $M = \frac{1}{2}(P_1 + P_2)$, and $D_{\text{KL}}(P \parallel Q) = \sum_c P(c) \log_2 \frac{P(c)}{Q(c)}$. Modalities with high $D_{\text{JS}}$ possess strong counterfactual decisiveness, guaranteeing diagnostic disambiguation regardless of which disease is true.

\subsection{Module 6: Multi-Factor Utility Scoring}
The optimal query $e^*$ is selected by maximizing the composite Net Policy Utility function $U(e)$:
\begin{equation}
    e^* = \operatorname*{arg\,max}_{e \in \mathcal{E}} \left[ \alpha \cdot D_{\text{JS}}(P_1 \parallel P_2) + \beta \cdot \mathbb{E}[\Delta H(e)] + \gamma \cdot \mathcal{S}(e) - \lambda \cdot C(e) - \mu \cdot R(e) \right]
\end{equation}
where $\mathbb{E}[\Delta H(e)]$ is predicted Expected Information Gain generated by the lightweight UtilityScoringNetwork MLP, $\mathcal{S}(e) = A[e, h_1, h_2]$ is the direct pairwise separation power, and $C(e)$ and $R(e)$ are physical acquisition cost and failure risk. The five Pareto hyperparameters are calibrated as: $\alpha = 1.0$, $\beta = 0.8$, $\gamma = 1.2$, $\lambda = 0.3$, and $\mu = 0.5$.

\subsection{Module 7: Optical Quality Gating (Safety Shield 1)}
To permanently eliminate Softmax hallucination on degraded imagery, Module 7 intercepts both initial photographs ($X_0$) and follow-up photographic clues prior to neural network ingestion:
\begin{enumerate}
    \item \textbf{Laplacian Blur Assessment:} Image sharpness is computed via the variance of the 2D discrete Laplacian convolution:
    \begin{equation}
        \sigma^2(I) = \frac{1}{N} \sum_{x,y} \left( \nabla^2 I_{\text{gray}}(x,y) - \mu_{\nabla} \right)^2
    \end{equation}
    where $\nabla^2 I = \frac{\partial^2 I}{\partial x^2} + \frac{\partial^2 I}{\partial y^2}$.
    \item \textbf{Luminance Assessment:} Mean illumination is computed as $\mu_L = \frac{1}{N} \sum_{x,y} I_{\text{gray}}(x,y)$.
\end{enumerate}

A composite optical quality score $q \in [0, 1]$ is generated:
\begin{equation}
    q = \min\left(1.0, \frac{\sigma^2(I)}{1.5 \cdot \tau_{\text{blur}}}\right) \cdot f_{\text{bright}}(\mu_L)
\end{equation}
where $\tau_{\text{blur}} = 80.0$\,$\sigma^2$, and $f_{\text{bright}}(\mu_L) = 1.0$ for $\mu_L \in [30, 230]$\,cd/m$^2$ (ramping linearly to 0 outside).
If $\sigma^2(I) < 80.0$ or $\mu_L \notin [30, 230]$, the image is marked REJECT ($q < 0.35$). The image is quarantined, model forward passes are bypassed, and the user is prompted to retake the photo without emitting erroneous disease probabilities.

\subsection{Module 8: Contradiction Verification and Specimen Mismatch Guard (Safety Shield 2)}
Module 8 protects the Bayesian posterior from corrupted or cross-specimen contaminant uploads through two distinct verification routines:
\begin{enumerate}
    \item \textbf{Directional Trajectory Contradiction:}
    Let $\Delta_{\text{exp}, 1} = P(h \mid e, h_1) - P_{\text{prior}}$ and $\Delta_{\text{exp}, 2} = P(h \mid e, h_2) - P_{\text{prior}}$ denote the planned counterfactual shift vectors, and $\Delta_{\text{obs}} = P_{\text{obs}} - P_{\text{prior}}$ denote the observed belief shift. Directional cosine alignment is:
    \begin{equation}
        \cos_k = \frac{\Delta_{\text{obs}} \cdot \Delta_{\text{exp}, k}}{\|\Delta_{\text{obs}}\| \|\Delta_{\text{exp}, k}\|}, \quad S_{\text{contra}} = \frac{1 - \max(\cos_1, \cos_2)}{2}
    \end{equation}
    If $S_{\text{contra}} \ge 0.40$, the observation triggers a confidence dampening penalty:
    \begin{equation}
        \rho = 1.0 - 0.50 \cdot S_{\text{contra}}
    \end{equation}
    \item \textbf{Specimen Mismatch Guard:}
    If incoming evidence exhibits high likelihood for an alien pathology $\theta_{\text{ev}}$ that was excluded during initial observation:
    \begin{equation}
        \theta_{\text{ev}} \notin \{h_1, h_2\} \;\wedge\; P_{\text{prior}}(\theta_{\text{ev}}) < 0.05 \;\wedge\; P_{\text{ev}}(\theta_{\text{ev}}) \ge 0.35
    \end{equation}
    the system detects a Specimen Mismatch (e.g., uploading a healthy leaf into an active blight inquiry). The evidence is completely QUARANTINED ($\rho = 0.0$), halting posterior corruption while retaining prior clinical beliefs.
\end{enumerate}

\subsection{Module 9: Reliability-Weighted Bayesian Belief Updater}
Valid evidence updates the diagnostic posterior according to the modified Bayesian update rule:
\begin{equation}
    P_{\text{new}}(\theta_c \mid e) = \frac{P_{\text{prior}}(\theta_c) \cdot \left[ L(e \mid \theta_c) \right]^\gamma}{\sum_{j=1}^K P_{\text{prior}}(\theta_j) \cdot \left[ L(e \mid \theta_j) \right]^\gamma}
\end{equation}
where $L(e \mid \theta_c)$ is the multimodal evidence likelihood, and the effective evidence exponent $\gamma$ dynamically scales with optical quality $q$ and contradiction factor $\rho$:
\begin{equation}
    \gamma = \operatorname{clip}(q \cdot \rho, 0.05, 1.5)
\end{equation}
On clean congruent evidence ($q=1.0, \rho=1.0$), $\gamma = 1.0$ (exact Bayesian update). On degraded inputs ($q < 0.7$), $\gamma$ attenuates belief shifts. On quarantined mismatches ($\rho = 0.0$), $\gamma \to 0.05$, neutralizing the contaminant.

\section{System Implementation and Hardware Architecture}

EvidenceGain-X is implemented in Python 3.10+ using PyTorch 2.2+, torchvision 0.17+, OpenCV-headless 4.8+, SciPy 1.11+, and Streamlit 1.30+. The architecture strictly decouples algorithmic modules into isolated subpackages (`src/models/`, `src/planning/`, `src/verification/`, `src/diagnosis/`).

To project textual field observations into the visual latent space, we designed a lightweight `TextEvidenceEncoder`. It maps a curated 35-term agricultural vocabulary covering lesion morphology (\textit{concentric}, \textit{rings}, \textit{halo}, \textit{pycnidia}), texture (\textit{water-soaked}, \textit{velvety}), and weather (\textit{humid}, \textit{heavy dew}) into a 64-dimensional embedding, followed by a projection head: $\text{Linear}(64, 128) \to \text{ReLU} \to \text{Linear}(128, 512) \to \text{LayerNorm}(512)$. Out-of-vocabulary tokens are hashed deterministically. LayerNorm aligns text embedding norms with ResNet-18 visual feature vectors ($\|\mathbf{h}\|_2$), enabling stable cosine distance comparisons.

\section{Experimental Setup}

\subsection{Dataset and Leakage-Free Partitioning}
Experiments were conducted using 200 authentic multi-modal episodes derived from the PlantVillage tomato corpus \cite{hughes2015}. To eliminate data leakage \cite{ferentinos2018}, partitioning was performed by physical plant identity:
\begin{itemize}
    \item \textbf{Training Set (70\%):} 140 episodes (strictly plant-isolated).
    \item \textbf{Validation Set (15\%):} 30 episodes (used for Platt calibration and hyperparameter tuning).
    \item \textbf{Testing Benchmark (15\%):} 24 fully validated multi-view episodes across the four target pathologies.
\end{itemize}
Zero physical plant overlap exists across partitions, ensuring true out-of-sample generalization.

\subsection{Evaluated Baseline Diagnostic Policies}
We benchmark EvidenceGain-X against six competitive decision policies:
\begin{itemize}
    \item \textbf{B0 (Initial Only):} Conventional passive single-shot ResNet-18.
    \item \textbf{B1 (Random Policy):} Acquires follow-up observations uniformly at random.
    \item \textbf{B2 (Fixed Order Sequence):} Rigid deterministic sequence ($e_1 \to e_2 \to e_3 \to e_4 \to e_5$).
    \item \textbf{B3 (Confidence-Driven):} Queries randomly until maximum probability exceeds 0.75.
    \item \textbf{B4 (Uncertainty-Only):} Queries modalities with highest generic historical entropy reduction without considering the rival pair.
    \item \textbf{B5 (EvidenceGain - Ablated):} Full counterfactual utility optimization without Module 7 \& 8 safety shields.
    \item \textbf{B6 (EvidenceGain-X - Proposed):} Complete nine-module framework with dual safety shields.
\end{itemize}

\section{Results and Quantitative Analysis}

\subsection{Diagnostic Accuracy and Efficiency}
Table~\ref{tab:benchmarks} summarizes quantitative performance across all seven policies on the 24-case benchmark.

\begin{table*}[t]
\caption{Quantitative Benchmark Evaluation Across 7 Decision Policies (24 Plant-Grouped Episodes)}
\label{tab:benchmarks}
\centering
\begin{tabular}{llccccc}
\toprule
\textbf{Policy Index} & \textbf{Diagnostic Policy Strategy} & \textbf{Accuracy (\%)} & \textbf{Macro F1} & \textbf{Mean Steps} & \textbf{$\Delta H$ (bits)} & \textbf{Adversarial Robustness} \\
\midrule
B0 & Initial Only (Passive Single-Shot) & 87.50\% & 0.8773 & 0.00 & 0.000 & 0.0\% (Hallucinates) \\
B1 & Random Acquisition & 87.50\% & 0.8726 & 1.00 & 0.221 & 20.0\% (Vulnerable) \\
B2 & Fixed Order Sequence & 91.67\% & 0.9188 & 1.00 & 0.284 & 30.0\% (Rigid) \\
B3 & Confidence-Driven ($P > 0.75$) & 87.50\% & 0.8726 & 0.50 & 0.203 & 15.0\% (Overconfident) \\
B4 & Uncertainty-Only (Generic EIG) & 83.33\% & 0.8233 & 1.12 & 0.229 & 25.0\% (Drifts) \\
B5 & EvidenceGain (Ablated w/o Safety) & 91.67\% & 0.9188 & 0.79 & 0.359 & 40.0\% (Poisoned) \\
\textbf{B6} & \textbf{EvidenceGain-X (Full Proposed System)} & \textbf{91.67\%} & \textbf{0.9188} & \textbf{0.88} & \textbf{0.359} & \textbf{100.0\% (Gated / Certified)} \\
\bottomrule
\end{tabular}
\end{table*}

Passive single-shot classification (B0) achieves 87.50\% top-1 accuracy and a Macro F1 of 0.8773, with all three misdiagnoses concentrated along the Early Blight vs. Septoria look-alike boundary. Random acquisition (B1) fails to improve accuracy (87.50\%) despite consuming 1.0 additional steps, proving that undirected querying introduces variance without resolving ambiguity. B4 degrades accuracy to 83.33\% by querying weather context when morphological lesion inspection was required.

Both B5 and EvidenceGain-X (B6) achieve peak diagnostic accuracy of \textbf{91.67\%} and a \textbf{Macro F1 of 0.9188}, resolving ambiguous look-alikes. EvidenceGain-X achieves this in an average of \textbf{0.88 steps} (a 12\% inquiry reduction compared to static acquisition B2), driven by conservative stopping on clean cases.

\subsection{Entropy Reduction Analysis}
EvidenceGain-X achieves an average entropy reduction of \textbf{0.359 bits} on ambiguous cases---a \textbf{62.4\% greater information gain} than random querying (0.221 bits) and 26.4\% higher than fixed sequencing (0.284 bits). This confirms that counterfactual simulation accurately identifies the single most discriminative physical clue for each specific confusion pair.

\subsection{Adversarial Stress-Testing and Safety Gating}
To evaluate cyber-physical resilience, 50 degraded and contradictory inputs were injected into the test pipeline:
\begin{itemize}
    \item \textbf{Defocus \& Motion Blur ($\sigma^2 < 80.0$):} B0 assigned an average 89.5\% confidence of Late Blight to meaningless blur. Module 7 intercepted \textbf{100\% of blurry uploads}, completely preventing Softmax hallucination.
    \item \textbf{Specimen Mismatch \& Contradiction:} When a healthy leaf was uploaded into an active Early Blight episode (where Healthy prior was 0.3\%), Module 8 detected specimen incongruence ($S_{\text{contra}} = 0.982$, prior $< 0.05$), quarantining the input and preserving prior diagnostic belief at 71.4\%. Baselines B0--B5 incorporated the foreign input, destabilizing the diagnosis.
\end{itemize}

\subsection{Component Ablation Analysis}
Table~\ref{tab:ablation} quantifies the operational necessity of each component.

\begin{table}[h]
\caption{System Component Ablation Study}
\label{tab:ablation}
\centering
\begin{tabular}{lccc}
\toprule
\textbf{Ablation Configuration} & \textbf{Accuracy} & \textbf{Mean Steps} & \textbf{Adversarial Pass} \\
\midrule
Full EvidenceGain-X (Proposed) & \textbf{91.67\%} & \textbf{0.88} & \textbf{100.0\%} \\
w/o Hypothesis Separation ($\alpha = 0$) & 87.50\% & 0.88 & 100.0\% \\
w/o Diagnostic Power ($\gamma = 0$) & 87.50\% & 0.92 & 100.0\% \\
w/o Uncertainty Reduction ($\beta = 0$) & 91.67\% & 1.12 & 100.0\% \\
w/o Cost Penalty ($\lambda = 0$) & 91.67\% & 1.25 & 100.0\% \\
w/o Optical Blur Gate (Module 7) & 58.33\% & 0.88 & 0.0\% \\
w/o Contradiction Guard (Module 8) & 75.00\% & 0.88 & 0.0\% \\
\bottomrule
\end{tabular}
\end{table}

Removing Module 7 causes catastrophic accuracy collapse to 58.33\% under field noise, proving that input quality gating is an indispensable prerequisite for trustworthy deployment. Removing cost penalization ($\lambda = 0$) increases farmer inquiry burden by 42\% (1.25 steps) without diagnostic gain.

\subsection{Latency and Edge Footprint}
Benchmarked on an Intel Core i7-12700K (single-thread CPU mode):
\begin{itemize}
    \item ResNet-18 forward pass: \textbf{12.1\,ms}
    \item Platt scaling \& entropy: \textbf{0.02\,ms}
    \item Module 7 Laplacian variance: \textbf{1.8\,ms}
    \item Module 5 Counterfactual simulation: \textbf{0.8\,ms}
    \item Module 8 Contradiction check: \textbf{0.4\,ms}
    \item \textbf{Total Diagnostic Loop Latency:} \textbf{$<$ 30\,ms} on CPU.
\end{itemize}
Memory footprint totals \textbf{44.5\,MB}, permitting offline deployment on sub-\$100 agricultural edge microcomputers.

\section{Limitations and Ethical Implications}
While EvidenceGain-X addresses key vulnerabilities of single-shot models, several limitations remain. First, the 24-episode benchmark, while rigorously plant-grouped, is limited in statistical scale due to the physical scarcity of complete multi-modal foliar episodes. Second, expanding to wider crop taxonomies (e.g., cassava, rice, maize) requires recalibrating the discriminative affinity matrix $\mathbf{A}$ through agronomic literature or automated knowledge-graph extraction. Ethically, EvidenceGain-X is designed as an advisory clinical decision-support tool; high-consequence interventions (such as chemical spraying or quarantine culling) should remain subject to human verification.

\section{Conclusion and Future Work}
This paper introduced \textbf{EvidenceGain-X}, an autonomous active diagnostic framework that resolves the core failure modes of single-shot agricultural AI. By combining temperature-scaled uncertainty quantification, counterfactual evidence planning, multi-objective utility optimization, and dual cyber-physical safety shields (optical blur gating and specimen mismatch detection), EvidenceGain-X elevates tomato pathology accuracy from 87.50\% to 91.67\% in an average of 0.88 sequential observation steps while achieving 100\% rejection against corrupted and contradictory inputs.

Future research will focus on INT8 TensorRT quantization for ultra-low-power microcontrollers, expanding the multimodal taxonomy to multi-crop disease complexes, and integrating conversational audio interfaces in regional languages to maximize accessibility for smallholder farming communities worldwide.

\begin{thebibliography}{00}
\bibitem{hughes2015} D. P. Hughes and M. Salath\'e, ``An open access repository of images on plant health to enable the development of mobile disease diagnostics,'' \textit{arXiv preprint arXiv:1511.08060}, 2015.
\bibitem{mohanty2016} S. P. Mohanty, D. P. Hughes, and M. Salath\'e, ``Using deep learning for image-based plant disease detection,'' \textit{Frontiers in Plant Science}, vol. 7, p. 1419, Sep. 2016.
\bibitem{agrios2005} G. N. Agrios, \textit{Plant Pathology}, 5th ed. Burlington, MA, USA: Elsevier Academic Press, 2005.
\bibitem{ferentinos2018} K. P. Ferentinos, ``Deep learning models for plant disease detection and diagnosis,'' \textit{Computers and Electronics in Agriculture}, vol. 145, pp. 311--318, Feb. 2018.
\bibitem{he2016} K. He, X. Zhang, S. Ren, and J. Sun, ``Deep residual learning for image recognition,'' in \textit{Proc. IEEE Conf. Comput. Vis. Pattern Recognit. (CVPR)}, Las Vegas, NV, USA, 2016, pp. 770--778.
\bibitem{ramcharan2017} A. Ramcharan, K. Baranowski, P. McCloskey, B. Ahmed, J. Legg, and D. P. Hughes, ``Deep learning for image-based cassava disease detection,'' \textit{Frontiers in Plant Science}, vol. 8, p. 1852, Oct. 2017.
\bibitem{guo2017} C. Guo, G. Pleiss, Y. Sun, and K. Q. Weinberger, ``On calibration of modern neural networks,'' in \textit{Proc. 34th Int. Conf. Mach. Learn. (ICML)}, Sydney, NSW, Australia, 2017, pp. 1321--1330.
\bibitem{gal2016} Y. Gal and Z. Ghahramani, ``Dropout as a Bayesian approximation: Representing model uncertainty in deep learning,'' in \textit{Proc. 33rd Int. Conf. Mach. Learn. (ICML)}, New York, NY, USA, 2016, pp. 1050--1059.
\bibitem{lakshminarayanan2017} B. Lakshminarayanan, A. Pritzel, and C. Blundell, ``Simple and scalable predictive uncertainty estimation using deep ensembles,'' in \textit{Adv. Neural Inf. Process. Syst. (NeurIPS)}, Long Beach, CA, USA, 2017, pp. 6402--6413.
\bibitem{settles2009} B. Settles, ``Active learning literature survey,'' Comput. Sci. Dept., Univ. Wisconsin-Madison, Tech. Rep. 1648, 2009.
\bibitem{chaloner1995} K. Chaloner and I. Verdinelli, ``Bayesian experimental design: A review,'' \textit{Stat. Sci.}, vol. 10, no. 3, pp. 273--304, Aug. 1995.
\bibitem{lindley1956} D. V. Lindley, ``On a measure of the information provided by an experiment,'' \textit{Ann. Math. Stat.}, vol. 27, no. 4, pp. 986--1005, Dec. 1956.
\bibitem{deng2009} J. Deng, W. Dong, R. Socher, L.-J. Li, K. Li, and L. Fei-Fei, ``ImageNet: A large-scale hierarchical image database,'' in \textit{Proc. IEEE Conf. Comput. Vis. Pattern Recognit. (CVPR)}, Miami, FL, USA, 2009, pp. 248--255.
\bibitem{pechmay2020} J. Pech-May, F. J. Sanchez-Marin, and J. A. Ramirez-Cortes, ``Laplacian-based focus measure operators for shape from focus: A comparative study,'' \textit{Opt. Lasers Eng.}, vol. 128, p. 106020, May 2020.
\bibitem{pearl2009} J. Pearl, \textit{Causality: Models, Reasoning, and Inference}, 2nd ed. Cambridge, UK: Cambridge Univ. Press, 2009.
\end{thebibliography}

\end{document}
